% !TeX root = ../../tesis.tex
% Esquema ilustrativo (decisión del autor 2026-11-04): claves y rutas
% tomadas de VFTModel/src/api/main.py; valores = marcadores de tipo.

\section{Ejemplo de Respuesta: Cobertura Espacial \texorpdfstring{($C$)}{(C)}}
\label{anx:vft-out-cobertura}

Esquema ilustrativo de la respuesta del indicador $C$
(§\ref{subsec:indicador-cobertura}). Los nombres de las claves y la ruta son los
de \texttt{VFTModel}; los valores son marcadores de tipo. Las cifras de la
corrida 2026-10-04 se reportan en el Capítulo~\ref{ch:analisis-red-actual}.

La respuesta trae la cobertura de cada demarcación y un resumen por dominio:
las entidades consultadas y la \zmvm{} de 76 demarcaciones.

\begin{lstlisting}[style=estiloPython,
  caption={Esquema de respuesta del endpoint \texttt{/spatial-coverage}
           (ilustrativo; \texttt{radio\_m} = 800).},
  label=lst:out-cobertura]
{
  "status": "success",
  "data": [
    { "Demarcacion": "<texto>", "Area_Total_km2": <número>,
      "Area_Cubierta_km2": <número>, "Cobertura_Porcentaje": <número> }
  ],
  "resumen_dominios": {
    "entidades": { "valor_bruto": <número>, "valor_normalizado": <número>,
                   "banda": "<texto>", "n_demarcaciones": <número>,
                   "area_total_km2": <número> },
    "zmvm_76":   { "valor_bruto": <número>, "valor_normalizado": <número>,
                   "banda": "<texto>", "n_demarcaciones": <número>,
                   "area_total_km2": <número> }
  }
}
\end{lstlisting}
